\section{如何读实验：Top-1、Top-2、步数、时间与 Scaling Law}

架构和系统约束明确后，实验的关键不是寻找一张“Switch 胜利图”，而是按顺序区分 quality per step、quality per wall-clock、固定 FLOPs 下的参数收益，以及通信带来的边际递减。讲座中的部分纵轴是 negative log perplexity；普通 perplexity（PPL）是交叉熵的指数，越小越好，而 negative log PPL 越大越好。

若 token 平均交叉熵为 $H$，则：

\[
\operatorname{PPL}=e^{H},\qquad -\log(\operatorname{PPL})=-H.
\]

其中，$H$ 是每 token 平均负对数似然；PPL 可粗略理解为模型面对的“有效选择数”，但它不直接等价于事实性、推理能力或下游任务质量；$-\log(\operatorname{PPL})$ 只是符号翻转，因此越大代表 loss 越低。

\subsection{Top-1 与 Top-2：容量因子改变 Pareto 前沿}

Top-2 在高 capacity 下每个 token 使用两个 experts，单步质量可能更好；但它也增加 FFN 计算和通信。当 capacity factor 降低时，Switch top-1 能以更低成本达到目标质量。正确结论不是“top-1 数学上更强”，而是它在讲座的硬件与容量设置中形成更好的 time-to-quality Pareto 点。

\lecturefigure{slide-17-comparison-moe-switch.jpg}{MoE top-2 与 Switch top-1 在不同 capacity factor 下的质量、速度与 time-to-quality。}{Stanford Online 官方视频 00:27:08--00:28:57。}

\begin{knowledgebox}{读图：比较顺序决定结论}
先看 capacity factor 是否相同，再看每步质量；随后看 speed 与达到同一质量阈值的时间。若 top-2 在 step 轴更好、top-1 在 time 轴更快，两者并不矛盾。最后检查低 capacity 是否造成更多 dropping，以及作者是否把省下的计算重新分配到 dense 层。
\end{knowledgebox}

\teachervoice{讲者将结论表述为 Pareto efficiency：高 capacity 的 top-2 可能有更好 step quality，但低 capacity top-1 减少通信、计算和内存，因此更快到达质量阈值。这个限定比一句“Switch beats MoE”准确得多。}

\begin{warningbox}{不要跨硬件照搬 Pareto 结论}
不同互连、batch size、expert placement、kernel 与编译器会改变通信占比。论文中的 time-to-quality 是特定系统的测量；迁移到 GPU 集群或低吞吐服务时必须重新 benchmark。
\end{warningbox}

\subsection{增加 experts：先看训练步，再看真实时间}

本节把同一组扩展实验放到两个横轴上阅读。固定每 token 只激活少量 experts 时，增加 expert 数会增加总参数与可分配容量。按训练步看，更多 experts 通常更快降低 loss，并获得更好渐近值；但曲线逐渐靠近，说明参数扩展存在边际递减。

\lecturefigure{slide-18-scaling-experts-per-training-step.jpg}{按训练步比较：expert 数增加带来质量提升，但收益逐渐递减。}{Stanford Online 官方视频 00:28:58--00:29:43。}

\begin{knowledgebox}{读图：训练步不是统一成本单位}
横轴每一步可能包含不同通信与 padding 开销。曲线向左或向上说明优化步效率更高，但不能直接推出训练更快。还应关注不同 expert 数下是否保持相同 batch、active FLOPs、capacity factor 与 sparse layer 频率。
\end{knowledgebox}

把横轴换成 wall-clock 后，通信成本进入图中。若更多 experts 跨越更多设备，dispatch 与 combine 会拖慢每一步；因此“每步收益”必须足以覆盖“每步更慢”，才会变成真实时间收益。

\lecturefigure{slide-19-scaling-experts-per-time.jpg}{按 wall-clock 比较：通信成本进入后，专家数的收益重新排序。}{Stanford Online 官方视频 00:29:44--00:30:51。}

\begin{importantbox}{实验报告的最小二联图}
至少同时报告 quality-versus-step 与 quality-versus-time。前者诊断优化和样本效率，后者回答系统是否真的更快。若再加 quality-versus-FLOPs 与成本，就能区分算法收益、硬件收益和资源转移。
\end{importantbox}

\subsection{Sparse scaling law：固定 FLOPs 增加参数仍有收益}

讲座进一步将 dense scaling law 扩展到 sparse setting：在每 token FLOPs 近似固定时，增加 sparsely activated parameters 仍可改善 loss。图中曲线不是无限线性下降；随着 experts 增加，收益逐渐变小，说明可用容量、数据覆盖、路由学习和通信共同限制扩展。

\lecturefigure{slide-20-sparse-scaling-laws.jpg}{固定每 token FLOPs 时，增加稀疏参数的经验 scaling 曲线。}{Stanford Online 官方视频 00:30:52--00:31:29。}

\begin{knowledgebox}{读图：固定 FLOPs 不等于固定总成本}
纵轴改善支持“条件参数有独立价值”，但横轴未完全包含参数存储、checkpoint、optimizer state、设备数量和通信。该图证明的是实验设置中的预测 loss 趋势，不证明无限加 experts 都具成本效益。
\end{knowledgebox}

\subsection{Expert parallelism、model parallelism 与小规模收益}

Expert parallelism 将不同 experts 放到不同设备；model parallelism 则切分同一 dense 权重或张量维度。二者可以互补：前者扩展条件参数池，后者让单个 dense block 也能跨设备。图中对比说明，组合后仍可继续改善，而不是只能二选一。

这里的 sharding（分片）是指把参数或张量沿某个维度切到不同设备，使每张设备只保存或计算其中一部分，而不再完整复制全部权重。Expert parallelism 可以视为以 expert 为天然分片单元的一种特殊布局。

\lecturefigure{slide-21-expert-vs-model-parallelism.jpg}{Expert parallelism 与 model parallelism 的质量对比及互补关系。}{Stanford Online 官方视频 00:31:30--00:32:57。}

\begin{knowledgebox}{概念解释：parallelism 不是一个词}
\textbf{Data parallelism} 复制模型、切分样本并同步梯度；\textbf{model parallelism} 切分同一层的参数或张量；\textbf{expert parallelism} 让不同设备拥有不同 experts，并通过 token all-to-all 调度。三者解决的内存和计算瓶颈不同。
\end{knowledgebox}

稀疏性也不只在超大规模才有效。讲座展示较小模型上仍有一致收益，这意味着研究者可以先在可承受规模验证 router、capacity 与系统实现，而不必把“训练万亿模型”作为使用 MoE 的先决条件。

\lecturefigure{slide-22-effective-at-small-scale.jpg}{Switch Transformer 在较小规模与不同 expert 数下仍显示收益。}{Stanford Online 官方视频 00:32:58--00:34:03。}

\teachervoice{老师强调小规模结果的意义：如果方法只在极端规模有效，就很难迭代和复现；在较小模型上观察一致趋势，允许团队先验证实现与训练稳定，再扩大设备数。}

最后一张曲线提醒：固定计算预算不断增加参数，最终会出现明显边际递减。原因可能包括每个 expert 接收的有效训练样本变少、router 难以学习可靠分区、参数存储与通信上升，以及数据本身不足以训练更多自由度。

\lecturefigure{slide-23-parameters-fixed-compute-diminishing-returns.jpg}{固定计算预算只增加参数，最终出现边际收益递减。}{Stanford Online 官方视频 00:34:04--00:36:25。}

\begin{warningbox}{“更多参数更样本高效”也有条件}
总数据固定时，专家越多，每个 expert 实际看到的 token 可能越少。若 router 分配极不均、某些专家长期冷启动，额外参数会变成未训练容量。应报告 expert utilization、每 expert token 数与长尾使用分布。
\end{warningbox}

\subsection{本章小结}

读 MoE 实验必须区分 step、time、FLOPs 与总系统成本。Top-1 的优势来自特定 capacity 和硬件上的 Pareto 权衡；更多 experts 提升质量，但通信、数据覆盖与路由学习带来边际递减。

